\documentclass[UTF8]{ctexart}
\usepackage{amsmath, amssymb, geometry}
\usepackage{xcolor}
\usepackage{enumitem}
\usepackage{tcolorbox}
\usepackage{cases}
\usepackage{fancyhdr}
\geometry{a4paper, margin=1in}

% 定义红色环境
\newenvironment{redenv}{\color{red}}{}

% 定义详细解析环境
\newenvironment{detailedanalysis}{%
    \color{red}
    \begin{enumerate}[label=\textbf{第\arabic*步：}, leftmargin=*, align=left, itemsep=1.5ex]
}{%
    \end{enumerate}
}

% 定义概念框
\newenvironment{conceptbox}{%
    \begin{tcolorbox}[colback=red!5, colframe=red!50!black, title=概念解析, fonttitle=\bfseries]
}{%
    \end{tcolorbox}
}

% 定义公式推导环境
\newenvironment{formuladerivation}{%
    \begin{enumerate}[label={\textbf{公式\arabic*：}}, leftmargin=*, align=left]
}{%
    \end{enumerate}
}

% 定义题目和解析的分隔线
\newcommand{\problemrule}{\noindent\rule{\textwidth}{1.5pt}\vspace{-0.8em}}
\newcommand{\analysisrule}{\noindent\rule{\textwidth}{0.8pt}\vspace{-0.8em}}

% 宏定义：方便排版选择题选项
\newcommand{\choicesFour}[4]{
    \begin{enumerate}[label=\Alph*., itemsep=0pt, parsep=0pt, topsep=5pt, labelsep=0.5em]
        \begin{minipage}{0.22\linewidth} \item #1 \end{minipage}
        \begin{minipage}{0.22\linewidth} \item #2 \end{minipage}
        \begin{minipage}{0.22\linewidth} \item #3 \end{minipage}
        \begin{minipage}{0.22\linewidth} \item #4 \end{minipage}
    \end{enumerate}
}
\newcommand{\choicesTwo}[4]{
    \begin{enumerate}[label=\Alph*., itemsep=0pt, parsep=0pt, topsep=5pt, labelsep=0.5em]
        \begin{minipage}{0.45\linewidth} \item #1 \end{minipage}
        \begin{minipage}{0.45\linewidth} \item #2 \end{minipage}
        \begin{minipage}{0.45\linewidth} \item #3 \end{minipage}
        \begin{minipage}{0.45\linewidth} \item #4 \end{minipage}
    \end{enumerate}
}
\newcommand{\choices}[4]{
    \begin{enumerate}[label=\Alph*., itemsep=0pt, parsep=0pt, topsep=5pt, labelsep=0.5em]
        \item #1
        \item #2
        \item #3
        \item #4
    \end{enumerate}
}

\newcommand{\blank}[1]{\underline{\hspace{#1}}}

\begin{document}

\section*{第六章 不定积分（提高）}

% ========== 第1题 ==========
\problemrule
\section*{【题目1】原函数概念}
设函数 $ f(x) $ 的一个原函数为 $ \frac{1}{x} $ ，则 $ f(x)= $ \blank{1cm}。
\choicesFour{$ -\frac{1}{x^{2}} $}{$ \frac{2}{x^{3}} $}{$ \frac{1}{x} $}{$ \ln|x| $}

\begin{redenv}
\analysisrule
\subsection*{逐步解析}

\begin{detailedanalysis}
\item \textbf{问题本质分析}
已知原函数求导函数。若 $F(x)$ 是 $f(x)$ 的原函数，则 $f(x) = F'(x)$。

\item \textbf{详细求解过程}
$F(x) = \frac{1}{x} = x^{-1}$。
$$ f(x) = F'(x) = (x^{-1})' = -1 \cdot x^{-2} = -\frac{1}{x^2} $$

\item \textbf{最终答案}
\textbf{答案}：A
\end{detailedanalysis}
\end{redenv}

% ========== 第2题 ==========
\problemrule
\section*{【题目2】分部积分法}
若 $ e^{-x} $ 是 $ f(x) $ 的一个原函数，则 $ \int xf(x)dx = $ \blank{1cm}。
\choicesTwo
{$ e^{-x}(1-x)+C $}
{$ e^{-x}(x+1)+C $}
{$ e^{-x}(x-1)+C $}
{$ -e^{-x}(x+1)+C $}

\begin{redenv}
\analysisrule
\subsection*{逐步解析}

\begin{detailedanalysis}
\item \textbf{详细求解过程}
已知 $\int f(x) dx = e^{-x} + C$，即 $f(x) = (e^{-x})' = -e^{-x}$。
计算不定积分：
$$ \int x f(x) dx = \int x (-e^{-x}) dx = - \int x e^{-x} dx $$
利用分部积分法：$\int u dv = uv - \int v du$。
令 $u = x, dv = e^{-x}dx \implies v = -e^{-x}$。
$$ \int x e^{-x} dx = x(-e^{-x}) - \int (-e^{-x}) dx = -xe^{-x} + \int e^{-x} dx = -xe^{-x} - e^{-x} $$
所以：
$$ \int x f(x) dx = -(-xe^{-x} - e^{-x}) + C = xe^{-x} + e^{-x} + C = e^{-x}(x+1) + C $$

\item \textbf{最终答案}
\textbf{答案}：B
\end{detailedanalysis}
\end{redenv}

% ========== 第3题 ==========
\problemrule
\section*{【题目3】导数与积分的关系}
设 $ \int f(x)dx = \ln(x + \sqrt{1 + x^{2}}) + C, $ 则 $ f(x) = $ \blank{1cm}。
\choicesFour{$ \frac{x}{\sqrt{1+x^{2}}} $}{$ \frac{1}{\sqrt{1+x^{2}}} $}{$ \frac{1}{1+x} $}{$ \frac{1}{1+x^{2}} $}

\begin{redenv}
\analysisrule
\subsection*{逐步解析}

\begin{detailedanalysis}
\item \textbf{详细求解过程}
$$ f(x) = [\ln(x + \sqrt{1+x^2})]' $$
这是常见导数。
$$ (\ln(x+\sqrt{1+x^2}))' = \frac{1}{x+\sqrt{1+x^2}} \cdot (1 + \frac{2x}{2\sqrt{1+x^2}}) = \frac{1}{\sqrt{1+x^2}} $$

\item \textbf{最终答案}
\textbf{答案}：B
\end{detailedanalysis}
\end{redenv}

% ========== 第4题 ==========
\problemrule
\section*{【题目4】指数函数积分}
$ \int 2^{3x} dx = $ \blank{1cm}。
\choicesFour{$ \frac{2^{3x}}{3\ln 2} + C $}{$ \frac{1}{3}(\ln 2)2^{3x} + C $}{$ \frac{2^{x}}{3} + C $}{$ \frac{2^{3x}}{\ln 2} + C $}

\begin{redenv}
\analysisrule
\subsection*{逐步解析}

\begin{detailedanalysis}
\item \textbf{详细求解过程}
利用公式 $\int a^u du = \frac{a^u}{\ln a} + C$。
$$ \int 2^{3x} dx = \frac{1}{3} \int 2^{3x} d(3x) = \frac{1}{3} \cdot \frac{2^{3x}}{\ln 2} + C = \frac{2^{3x}}{3\ln 2} + C $$

\item \textbf{最终答案}
\textbf{答案}：A
\end{detailedanalysis}
\end{redenv}

% ========== 第5题 ==========
\problemrule
\section*{【题目5】变上限积分求导}
设函数 $ f(x) $ 在 $ [a, b] $ 上连续， $ F(x) = \int_{0}^{x^{3}} f(t) dt $ ，则 $ F'(x) = $ \blank{1cm}。
\choicesFour{$ f(x) $}{$ f(x^{3}) $}{$ 3x^{2}f(x^{3}) $}{$ 3x^{2}f(x) $}

\begin{redenv}
\analysisrule
\subsection*{逐步解析}

\begin{detailedanalysis}
\item \textbf{详细求解过程}
复合函数链式法则：
$$ \frac{d}{dx} \int_0^{g(x)} f(t) dt = f(g(x)) \cdot g'(x) $$
本题中 $g(x) = x^3$。
$$ F'(x) = f(x^3) \cdot (x^3)' = f(x^3) \cdot 3x^2 $$

\item \textbf{最终答案}
\textbf{答案}：C
\end{detailedanalysis}
\end{redenv}

% ========== 第6题 ==========
\problemrule
\section*{【题目6】凑微分法还原函数}
设 $ f'(x^{2}) = \frac{1}{x} (x > 0) $ ，则 $ f(x) = $ \blank{2cm}。

\begin{redenv}
\analysisrule
\subsection*{逐步解析}

\begin{detailedanalysis}
\item \textbf{详细求解过程}
\begin{enumerate}[label=(\arabic*)]
\item \textbf{换元}：令 $t = x^2$，则 $x = \sqrt{t}$（因 $x>0$）。
\item \textbf{代入}：
$f'(t) = \frac{1}{\sqrt{t}} = t^{-1/2}$。
\item \textbf{求原函数}：
$f(t) = \int t^{-1/2} dt = 2t^{1/2} + C = 2\sqrt{t} + C$。
$f(x) = 2\sqrt{x} + C$。
\end{enumerate}

\item \textbf{最终答案}
\textbf{答案}：$ 2\sqrt{x} + C $
\end{detailedanalysis}
\end{redenv}

% ========== 第7题 ==========
\problemrule
\section*{【题目7】分部积分}
设 $ f(x) $ 的一个原函数是 $ \sin x, $ 则 $ \int x f'(x)dx = $ \blank{2cm}。

\begin{redenv}
\analysisrule
\subsection*{逐步解析}

\begin{detailedanalysis}
\item \textbf{详细求解过程}
由题可知 $f(x) = (\sin x)' = \cos x$。
所以 $f'(x) = (\cos x)' = -\sin x$。
积分 $I = \int x (-\sin x) dx = - \int x \sin x dx$。
分部积分：
$$ \int x \sin x dx = x(-\cos x) - \int (-\cos x) dx = -x\cos x + \sin x + C $$
所以 $I = -(-x\cos x + \sin x) + C = x\cos x - \sin x + C$。

\item \textbf{最终答案}
\textbf{答案}：$ x\cos x - \sin x + C $
\end{detailedanalysis}
\end{redenv}

% ========== 第8题 ==========
\problemrule
\section*{【题目8】合并底数积分}
$ \int 3^{x}e^{x}dx = $ \blank{2cm}。

\begin{redenv}
\analysisrule
\subsection*{逐步解析}

\begin{detailedanalysis}
\item \textbf{详细求解过程}
$3^x e^x = (3e)^x$。
利用 $\int a^x dx = \frac{a^x}{\ln a} + C$。
$$ \int (3e)^x dx = \frac{(3e)^x}{\ln(3e)} + C = \frac{3^x e^x}{1 + \ln 3} + C $$

\item \textbf{最终答案}
\textbf{答案}：$ \frac{3^x e^x}{1+\ln 3} + C $
\end{detailedanalysis}
\end{redenv}

% ========== 第9题 ==========
\problemrule
\section*{【题目9】拆项积分}
不定积分 $ \int\frac{x-3}{x^{2}+1}dx= $ \blank{2cm}。

\begin{redenv}
\analysisrule
\subsection*{逐步解析}

\begin{detailedanalysis}
\item \textbf{详细求解过程}
拆分为两项：
$$ \int \frac{x}{x^2+1} dx - \int \frac{3}{x^2+1} dx $$
第一项（凑微分）：$\frac{1}{2} \int \frac{d(x^2+1)}{x^2+1} = \frac{1}{2} \ln(x^2+1)$。
第二项（基本公式）：$3 \arctan x$。
结果：$\frac{1}{2}\ln(1+x^2) - 3\arctan x + C$。

\item \textbf{最终答案}
\textbf{答案}：$ \frac{1}{2}\ln(1+x^2) - 3\arctan x + C $
\end{detailedanalysis}
\end{redenv}

% ========== 第10题 ==========
\problemrule
\section*{【题目10】辅助变换积分}
不定积分 $ \int \frac{1}{e^{x} + e^{-x}} dx = $ \blank{2cm}。

\begin{redenv}
\analysisrule
\subsection*{逐步解析}

\begin{detailedanalysis}
\item \textbf{详细求解过程}
分子分母同乘 $e^x$：
$$ \int \frac{e^x}{e^{2x} + 1} dx $$
令 $u = e^x$，则 $du = e^x dx$。
$$ \int \frac{du}{u^2+1} = \arctan u + C = \arctan(e^x) + C $$

\item \textbf{最终答案}
\textbf{答案}：$ \arctan e^x + C $
\end{detailedanalysis}
\end{redenv}

% ========== 第11题 ==========
\problemrule
\section*{【题目11】不定积分计算}
已知 $ f(x) $ 的一个原函数为 $ \frac{\sin x}{1 + x \sin x} $ ，求 $ \int f(x) f'(x) dx $。

\begin{redenv}
\analysisrule
\subsection*{逐步解析}

\begin{detailedanalysis}
\item \textbf{详细求解过程}
\begin{enumerate}[label=(\arabic*)]
\item \textbf{积分形式分析}：
$$ \int f(x) f'(x) dx = \int f(x) d(f(x)) = \frac{1}{2} f^2(x) + C $$
\item \textbf{求 f(x)}：
设 $F(x) = \frac{\sin x}{1 + x \sin x}$，则 $f(x) = F'(x)$。
题意是求积分结果，即 $\frac{1}{2}f^2(x)+C$。
需要算出 $f(x)$ 吗？根据答案习惯，如果能算最好算出来；如果式子太复杂，通常写成 $f(x)$ 的形式，但这里可能是直接代入。
$F'(x) = \frac{\cos x(1+x\sin x) - \sin x(1 \cdot \sin x + x\cos x)}{(1+x\sin x)^2} = \frac{\cos x + x\sin x \cos x - \sin^2 x - x\sin x \cos x}{(1+x\sin x)^2} = \frac{\cos x - \sin^2 x}{(1+x\sin x)^2}$。
所以结果是 $\frac{1}{2} \left[ \frac{\cos x - \sin^2 x}{(1+x\sin x)^2} \right]^2 + C$。
\textbf{注}：若是填空题，可能直接写 $\frac{1}{2}f^2(x)+C$ 且说明 $f(x)$ 是 $F(x)$ 导数即可，但最好写具体。由于式子复杂，不排除题目可能有误或有简便法，但直接套公式 $\frac{1}{2}f^2(x)$ 是最稳妥的。
\end{enumerate}

\item \textbf{最终答案}
\textbf{答案}：$ \frac{1}{2} \left[ \frac{\cos x - \sin^2 x}{(1+x\sin x)^2} \right]^2 + C $
\end{detailedanalysis}
\end{redenv}

% ========== 第12题 ==========
\problemrule
\section*{【题目12】求原函数表达式}
若 $ \int xf(x)dx = \arcsin x + C $ 求 $ I = \int \frac{1}{f(x)}dx $。

\begin{redenv}
\analysisrule
\subsection*{逐步解析}

\begin{detailedanalysis}
\item \textbf{详细求解过程}
\begin{enumerate}[label=(\arabic*)]
\item \textbf{求 f(x)}：
两边求导：$xf(x) = (\arcsin x)' = \frac{1}{\sqrt{1-x^2}}$。
所以 $f(x) = \frac{1}{x\sqrt{1-x^2}}$。
\item \textbf{计算积分}：
$$ \int \frac{1}{f(x)} dx = \int x\sqrt{1-x^2} dx $$
凑微分：$d(1-x^2) = -2x dx$。
$$ = -\frac{1}{2} \int (1-x^2)^{1/2} d(1-x^2) = -\frac{1}{2} \cdot \frac{2}{3} (1-x^2)^{3/2} + C $$
$$ = -\frac{1}{3} (1-x^2)^{3/2} + C $$
\end{enumerate}

\item \textbf{最终答案}
\textbf{答案}：$ -\frac{1}{3} (1-x^2)^{3/2} + C $
\end{detailedanalysis}
\end{redenv}

% ========== 第13题 ==========
\problemrule
\section*{【题目13】换元积分法}
求不定积分 $ \int\frac{e^{\sqrt{x}}\sin\sqrt{x}}{2\sqrt{x}}dx $。

\begin{redenv}
\analysisrule
\subsection*{逐步解析}

\begin{detailedanalysis}
\item \textbf{详细求解过程}
令 $t = \sqrt{x}$，则 $dt = \frac{1}{2\sqrt{x}} dx$。
$$ \int e^{\sqrt{x}}\sin\sqrt{x} \cdot \frac{1}{2\sqrt{x}} dx = \int e^t \sin t dt $$
利用两次分部积分公式：
$\int e^t \sin t dt = \frac{e^t(\sin t - \cos t)}{2} + C$。
回代 $t = \sqrt{x}$：
$$ \frac{e^{\sqrt{x}}(\sin\sqrt{x} - \cos\sqrt{x})}{2} + C $$

\item \textbf{最终答案}
\textbf{答案}：$ \frac{e^{\sqrt{x}}(\sin\sqrt{x} - \cos\sqrt{x})}{2} + C $
\end{detailedanalysis}
\end{redenv}

% ========== 第14题 ==========
\problemrule
\section*{【题目14】三角函数积分}
计算 $ \int \sin^{2} x \cos^{3} x dx $。

\begin{redenv}
\analysisrule
\subsection*{逐步解析}

\begin{detailedanalysis}
\item \textbf{详细求解过程}
留一个 $\cos x$ 凑微分，其余化为 $\sin x$。
$$ \int \sin^2 x \cos^2 x \cdot \cos x dx = \int \sin^2 x (1-\sin^2 x) d(\sin x) $$
令 $u = \sin x$：
$$ \int u^2 (1-u^2) du = \int (u^2 - u^4) du = \frac{u^3}{3} - \frac{u^5}{5} + C $$
回代：
$$ \frac{1}{3}\sin^3 x - \frac{1}{5}\sin^5 x + C $$

\item \textbf{最终答案}
\textbf{答案}：$ \frac{1}{3}\sin^3 x - \frac{1}{5}\sin^5 x + C $
\end{detailedanalysis}
\end{redenv}

% ========== 第15题 ==========
\problemrule
\section*{【题目15】凑微分法}
计算 $ \int \frac{\sqrt{\arcsin x}}{\sqrt{1-x^{2}}} dx $。

\begin{redenv}
\analysisrule
\subsection*{逐步解析}

\begin{detailedanalysis}
\item \textbf{详细求解过程}
观察到 $(\arcsin x)' = \frac{1}{\sqrt{1-x^2}}$。
凑微分：
$$ \int \sqrt{\arcsin x} d(\arcsin x) $$
令 $u = \arcsin x$：
$$ \int u^{1/2} du = \frac{2}{3} u^{3/2} + C = \frac{2}{3} (\arcsin x)^{3/2} + C $$

\item \textbf{最终答案}
\textbf{答案}：$ \frac{2}{3} (\arcsin x)^{3/2} + C $
\end{detailedanalysis}
\end{redenv}

\end{document}
